% from ega4-2-fr.tex
\oldpage[IV-2]{202}

donc (0$_{\mathrm I}$, 7.4.4) l'idéal $\mathfrak mV$ n'est pas un idéal de
définition de $V$. Comme $V$ est un anneau local intègre noethérien de
dimension 1, 0 est le seul idéal de $V$ qui ne soit pas un idéal de
définition, donc $\mathfrak mV=0$. Mais on a $A\subset V$ et $V$ est
intègre, d'où $\mathfrak m=0$ et $A=k$ est un corps. On en déduit tout
d'abord $\dim(C)\leqslant1$ (0, 16.1.5); mais comme $\dim(V)=1$, les
relations
\[
\dim(V)\leqslant\dim(B')=\dim(B)\leqslant\dim(C)
\]
montrent que cela entraîne
$\dim(C)=\dim(B)=\dim(B')=\dim(B'/\mathfrak p')=1$, et par suite
$\mathfrak p'$ est nécessairement un idéal premier minimal de $B'$. Nous
arriverons donc encore à une contradiction si nous prouvons que
$B'_{\mathfrak p'}$ est un corps, ou encore que l'anneau $B'$ est réduit.
Or, comme $C$ est une $k$-algèbre de type fini, la clôture intégrale $C_1$
de $C$ est une $C$-algèbre finie (Bourbaki, \emph{Alg. comm.}, chap. V,
\S~3, n\textsuperscript{o}~2, th. 2); si l'on pose
$S=C-\mathfrak n$, $B_1=S^{-1}C_1$ est la clôture intégrale de $B$, donc
$B_1$ est une $B$-algèbre finie, et par suite un anneau semi-local
noethérien, intègre et intégralement clos et de dimension 1 (0, 16.1.5); si
$\mathfrak m_j$ $(1\leqslant j\leqslant h)$ sont ses idéaux maximaux, les
$(B_1)_{\mathfrak m_j}$ sont donc des anneaux de valuation discrète (II,
7.1.6), et le complété $B'_1$ de $B_1$ est le composé direct des anneaux de
valuation discrète complétés des $(B_1)_{\mathfrak m_j}$ (Bourbaki,
\emph{Alg. comm.}, chap. III, \S~2, n\textsuperscript{o}~13, prop. 18);
$B'_1$ est donc réduit, et comme le complété $B'$ de $B$ est un
sous-anneau de $B'_1$ (Bourbaki, \emph{Alg. comm.}, chap. IV, \S~2,
n\textsuperscript{o}~5, cor. 3 de la prop. 9), c'est aussi un anneau réduit.
C.Q.F.D.
\end{proof}

\begin{corollary}[7.4.5]
\label{IV.2.s7.4.5-fr}
Supposons que la propriété $\boldsymbol P$ satisfasse aux conditions
$(\mathrm P_{\mathrm I})$, $(\mathrm P_{\mathrm{II}})$,
$(\mathrm P_{\mathrm{III}})$. Soit $A$ un anneau noethérien. Les conditions
suivantes sont équivalentes :
\begin{enumerate}
\item[\emph{a)}] Pour tout idéal premier $\mathfrak p$ de $A$,
$A_{\mathfrak p}$ est un $\boldsymbol P$-anneau.
\item[\emph{b)}] Pour tout idéal maximal $\mathfrak m$ de $A$,
$A_{\mathfrak m}$ est un $\boldsymbol P$-anneau.
\end{enumerate}
Si de plus $\boldsymbol P$ satisfait à $(\mathrm P_{\mathrm{IV}})$, les
propriétés a) et b) sont aussi équivalentes à :
\begin{enumerate}
\item[\emph{c)}] Pour toute $A$-algèbre de type fini $B$ et tout idéal
premier $\mathfrak q$ de $B$, $B_{\mathfrak q}$ est un
$\boldsymbol P$-anneau.
\end{enumerate}
\end{corollary}

\begin{proof}
L'équivalence de a) et b) résulte de (7.4.4, (ii)), et celle de a) et c) de
(7.4.4, (ii)).

Lorsque la condition a) de (7.4.5) est remplie, on dit que $A$ est un
$\boldsymbol P$-anneau; pour les anneaux semi-locaux noethériens, cette
définition coïncide avec celle de (7.3.13), lorsque les conditions
$(\mathrm P_{\mathrm I})$, $(\mathrm P_{\mathrm{II}})$ et
$(\mathrm P_{\mathrm{III}})$ sont satisfaites. L'anneau $\mathbf Z$ est un
$\boldsymbol P$-anneau (7.3.19, (iv)); tout anneau local complet est un
$\boldsymbol P$-anneau.
\end{proof}

\begin{proposition}[7.4.6]
\label{IV.2.s7.4.6-fr}
Supposons que la propriété $\boldsymbol P$ satisfasse aux conditions
$(\mathrm P_{\mathrm I})$, $(\mathrm P_{\mathrm{II}})$ et
$(\mathrm P_{\mathrm{III}})$. Soient $A$ un anneau noethérien,
$\mathfrak J$ un idéal de $A$, $\widehat A$ le séparé complété de $A$ pour
la topologie $\mathfrak J$-préadique. Alors, si $A$ est un
$\boldsymbol P$-anneau (7.4.5), le morphisme canonique
$\operatorname{Spec}(\widehat A)\to\operatorname{Spec}(A)$ est un
$\boldsymbol P$-morphisme.
\end{proposition}

\begin{proof}
En utilisant (7.3.2, c$'$)), il suffit de prouver que pour tout idéal
maximal $\mathfrak n$ de $\widehat A$, d'image réciproque $\mathfrak m$
dans $A$, le morphisme
$\operatorname{Spec}((\widehat A)_{\mathfrak n})\to
\operatorname{Spec}(A_{\mathfrak m})$ est un $\boldsymbol P$-morphisme.
On sait (Bourbaki, \emph{Alg. comm.}, chap. III, \S~3,
n\textsuperscript{o}~4, prop. 8) que l'homomorphisme canonique
$A_{\mathfrak m}\to(\widehat A)_{\mathfrak n}$ est injectif, que la
topologie $\mathfrak mA_{\mathfrak m}$-préadique sur $A_{\mathfrak m}$ est
induite par la topologie
$\mathfrak n(\widehat A)_{\mathfrak n}$-préadique et que
$A_{\mathfrak m}$ est dense dans $(\widehat A)_{\mathfrak n}$ pour cette
topologie, de sorte que le complété de $A_{\mathfrak m}$ pour la topologie
$\mathfrak mA_{\mathfrak m}$-préadique s'identifie à celui de
$(\widehat A)_{\mathfrak n}$ pour la topologie
$\mathfrak n(\widehat A)_{\mathfrak n}$-préadique. On a donc deux morphismes
\[
\operatorname{Spec}(\widehat{A_{\mathfrak m}})
\xrightarrow{f}\operatorname{Spec}((\widehat A)_{\mathfrak n})
\xrightarrow{g}\operatorname{Spec}(A_{\mathfrak m})
\]

